\section{The audited closure architecture}
\label{sec:arch}
% ============================================================================

This section defines the paper's central object. It is a \emph{contract}: a
declaration of what a structure must exhibit, under what audits, for the
catalog's results to compose. It is not an exhibited system, and
Section~\ref{sec:t17} keeps the distance between the contract and any
instantiation explicit.

\subsection{The full C-package}
\label{sec:cpackage}

\begin{definition}[Full C-package target]
\label{def:cpackage}
The \emph{full C-package} $\Cpkg$ is the eight-clause full-design and audit
contract $\mathrm{C1}$--$\mathrm{C8}$, together with cross-clause coherence
data and a full promotion-gate panel. A candidate consists of a carrier $X$,
declared transition/action families, a public lens and exposure channels, a
predictive self-lens, branch structure, a maintenance ledger, and protocol
generators; the clauses constrain that data as follows.
\begin{description}
\item[C1 (objecthood).] A declared multiscale exposure-closure family with
robust fixed-point structure: closure operators $E_\tau$ at declared stagings
$\tau$, exact fixed-point and defect bookkeeping, and perturbation audits with
declared retention thresholds.
\item[C2 (strict promotion).] A strict promoted distinction across declared
layers: $\Dfact(\pi_0,\pi_1)\neq\emptyset$ for the declared lower and promoted
readouts, certified by a split pair.
\item[C3 (reflexive self-coordinate).] A recursively stable self-coordinate:
refinement of the public self-lens by dynamical profiles terminates at a
predictive quotient that carries the package's own continuation coordinate.
\item[C4 (non-exposure).] The determining predictive state is not exposed
relative to the declared public channel family: no admissible current
observable factors through it non-constantly on the realized support.
Hiddenness is channel-relative by construction.
\item[C5 (drive/credit).] A declared audit of action credit: generator
provenance routes computed directionality into \emph{self}, \emph{other}, or
\emph{none}, with scheduled and internalized controls that must not be
credited as self.
\item[C6 (temporal closure).] A maintained, funded order/tick/arrow structure:
a staged phase carrier that persists and progresses under budget, with the
arrow carried by the package's own audit channel rather than by external
scheduling.
\item[C7 (binding).] Strict recombination refinement over the branchwise
quotient: joint content over declared branches, witnessed by
$\Rrec \succ \Kbw$ (Section~\ref{sec:structure}).
\item[C8 (viability).] Nonempty maintained viability: $\Kviab\neq\emptyset$
under the declared maintenance ledger, with underfunding collapsing the
kernel.
\end{description}
\end{definition}

The full design object also names four \emph{derived} structures that are
consequences of the clauses rather than clauses themselves: D1, the dual-price
structure of binding C8 constraints (Section~\ref{sec:structure}); D2, the
descended quality geometry of the package's own states
(Section~\ref{sec:structure}); D3, a proposed scheduling/marking reading of
attention-like allocation, which is recognition-grade only (the present
repository implements no price-based attention scheduler); and D4, the declared
report/descent channel with its no-overreading boundary
(Section~\ref{sec:interface}).

Two premises sit \emph{outside} the structural boundary and are
stated here only so the reader can see where they will be spent
(Section~\ref{sec:t17}): A1 (Ontic Closure Completeness) and A2 (Substrate
Indifference) are philosophical postulates, not clauses, and no structural
result in Sections~\ref{sec:t1}--\ref{sec:clusters} depends on them.

\begin{table}[t]
\centering\footnotesize
\begin{tabular}{p{0.06\linewidth}p{0.2\linewidth}p{0.24\linewidth}p{0.22\linewidth}p{0.18\linewidth}}
\toprule
Clause & Standard object & Strongest landed evidence & Simplified (T17-S) proxy & Blocking condition / nonclaim \\
\midrule
C1 & closure operators $E_\tau$ with fixed-point, defect, and perturbation audits & T1: exact $E_\tau$ audit on the ten-state kernel; unique perturbed fixed point retaining $77/80$ CORE mass at $\varepsilon=1/16$ & mechanics-coupled closure with staged-transition derivation and one-step audits & a fixed point is not an object or phenomenal entity without a recognition bridge \\
C2 & factorization defect $\Dfact\neq\emptyset$ & T1 split pair; T3 nonfactorization corollary & whole-carrier defect from strict ledger & no experiential novelty or universal explanatory gap \\
C3 & terminal profile refinement to predictive quotient & T1: kernel-restricted refinement $6\to7\to7$, terminal $=$ canonical $M$ on the kernel & one-step complete-observation-stable quotient & one-step stability at generic grade is not recursive terminality in full generality \\
C4 & non-exposure of determining state & T2: connector loop moves $M$ while fixing $Q$; dissolution control & single collision-generated class under the complete joint observation & no absolute privacy; expanded interfaces may expose \\
C5 & credit rule over generator provenance & T5: exact ledgers credit self/other/none; DPI null blocked & aggregate no-fake-arrow balance control & aggregate balance does not instantiate the concrete two-component lift \\
C6 & funded order/tick/arrow & T8: funded trace with monotone ledger, $\Delta A = 2$; unfunded stall & budget-disciplined trace proxy & the trace proxy is not the full imported clock theorem \\
C7 & strict recombination refinement & T7: $\Rrec\succ\Kbw$ witness; marking/erasure controls & global refinement with nontrivial fiber & no phenomenal-unity derivation \\
C8 & viability kernel under maintenance & T1 deadline contrast; T11 funded/underfunded contrast & strict viability shrinkage under collapse budget & shrinkage is proved generically, not universal empty collapse \\
\bottomrule
\end{tabular}
\caption{\textbf{Tbl-2: the C-package, compactly.} Each clause with its
standard mathematical object, its strongest landed component evidence, the
simplified certificate's proxy for it, and the inference it blocks.
Scope/price: component witnesses live on separately declared carriers; no
integrated instance exists (Section~\ref{sec:t17}).}
\label{tbl:cpackage}
\end{table}

\begin{figure}[t]
\centering
\fbox{\parbox{0.92\linewidth}{\small\emph{[Figure F2: C-package dependency
graph; rendered per the artifact figure plan.]} Nodes C1--C8 with dependency
edges; derived nodes D1--D4 attached to their source clauses; A1/A2 drawn
outside the structural boundary; the simplified certificate $\Certsimp$ shown
as a distinct inset, with no ordering arrow between it and the full package.}}
\caption{\textbf{F2: the C-package as a dependency graph.} Takeaway: the
clauses are one contract with declared interactions, not eight independent
slogans. Scope restriction: the graph shows the design target; it does not
assert an integrated instance, and no equivalence or ordering between the full
package and the simplified certificate is implied.}
\label{fig:cpackage}
\end{figure}

\subsection{Dependencies and independent axes}

The clauses are not symmetric. A forcing cascade organizes them: under
deadline and live-signal conditions, C3's self-coordinate is \emph{forced}
(Section~\ref{sec:t1}); C4 hiddenness, C6 temporal closure, and C2 strictness
are \emph{generic or declared}, holding under broad, explicitly stated
interface conditions; and C7 binding and C5 drive are \emph{contingent axes}
that a package may simply lack. The two contingent axes generate a four-cell
classifier lattice of degraded packages (T1b): binding and drive both absent
(a controller), binding without drive, drive without binding, and the full
subject cell. The cells are axis-independently witnessed classifications
(T7 witnesses each binding value and T5 each drive value on their own
carriers), not four jointly built systems; the lattice is a classificatory
schema, not a claim that these are the only possible structures.

The same synthesis packet lands a scoped recognition audit \gradeR{}: for this
catalog, every C1--C8 clause maps to composites of the declared six roles
without introducing a seventh irreducible role. The audit is scoped to this
package and the current role vocabulary; it does not prove any universal
completeness of the six roles.

\subsection{Why single ingredients are insufficient}
\label{sec:insufficient}

Each familiar single-ingredient proposal reappears inside the contract as one
clause, and each is separated from sufficiency by an explicit control
elsewhere in the catalog. Agency or regulation alone is C5/C8 material:
the credit audit exists precisely because apparent drive can be externally
scheduled or internalized without being creditable to self
(Section~\ref{sec:selfaudit}), and a nonempty viability kernel is controlled
persistence, not experience. Persistence alone (C1/C8) is exhibited by any
maintained fixed point; the T4 study shows a certificate-free reading of the
same public face (Section~\ref{sec:interface}). Integration or binding alone
(C7) collapses under marking while the branch family survives
(Section~\ref{sec:structure}). Report alone (D4) licenses only descended
boundary content, and the report map is provably non-injective on hidden
states (Section~\ref{sec:interface}). The contract is the conjunction plus its
coherence conditions; the catalog's controls are what make each ``alone''
claim fail exactly rather than rhetorically.

\subsection{Where bridges enter}

Three places, and only three, connect the structural contract to
non-structural readings. First, recognition bridges \gradeR{}: the named
proposals that D1's dual prices have the profile of felt valence
(Section~\ref{sec:structure}), that clause patterns correspond to
phenomenological categories, and that D3's marking dynamics point at
attention. Each is retractable without touching Tier~S. Second, the
ontological postulates A1/A2, spent only in Section~\ref{sec:t17}. Third, the
prospective empirical interface, which is specified as unmet obligations in
Section~\ref{sec:empirical}. Everything else in the paper is lane S.

\subsection{The full target versus the simplified certificate}
\label{sec:fullvssimp}

One distinction must be visible before any result is stated, because the
realization section depends on it. The full C-package $\Cpkg$ of
Definition~\ref{def:cpackage} is the design target. What is currently
\emph{mechanized} is a simplified, carrier-indexed certificate schema
$\Certsimp(X)$ (Section~\ref{sec:t17}): mechanics-based clause predicates
C1--C8, cross-clause coherence, and a simplified promotion ledger. The schemas
are not ordered by any landed theorem: the simplified certificate uses a
weaker imported gate ledger while adding one-step coherence restrictions
absent from the full design definition; no implication or equivalence between
them has been proved or is claimed. Three specific weakenings recur and are
flagged where used: the simplified C3 is one-step stable rather than
recursively terminal in full generality; the simplified C6 is a
budget-disciplined trace proxy rather than the full imported clock theorem;
and the simplified C8 proves strict viability \emph{shrinkage} under a
dominated collapse budget, not universal empty collapse. Table~\ref{tbl:cpackage}
carries these boundaries clause by clause.

% ============================================================================
